\documentclass[12pt]{article}

\usepackage{amsmath, amssymb, amsthm}
\usepackage{geometry}
\usepackage{hyperref}
\usepackage{enumitem}
\usepackage{xcolor}
\usepackage{mdframed}
\usepackage{booktabs}
\usepackage{graphicx}

\geometry{margin=1in}

\definecolor{keycolor}{RGB}{180, 70, 0}
\definecolor{notecolor}{RGB}{240, 240, 255}
\definecolor{takeawaycolor}{RGB}{255, 245, 230}

\newmdenv[
  backgroundcolor=notecolor,
  linecolor=keycolor,
  linewidth=1.5pt,
  topline=false,
  bottomline=false,
  rightline=false,
  leftmargin=0pt,
  rightmargin=0pt,
  innerleftmargin=10pt,
  innerrightmargin=10pt,
]{notebox}

\newmdenv[
  backgroundcolor=takeawaycolor,
  linecolor=keycolor,
  linewidth=2pt,
  roundcorner=5pt,
  innerleftmargin=10pt,
  innerrightmargin=10pt,
  innertopmargin=8pt,
  innerbottommargin=8pt,
]{takeawaybox}

\newtheorem{theorem}{Theorem}
\newtheorem{definition}{Definition}

\title{
  \textbf{CS 3600: Introduction to Artificial Intelligence}\\[0.5em]
  \large Lecture 10 --- Neural Networks\\[0.3em]
  \normalsize Aaron Hillegass, Georgia Tech
}
\author{Lecture Notes}
\date{}

\begin{document}

\maketitle
\tableofcontents
\newpage

% ---------------------------------------------------------------
\section{Motivation: Why Neural Networks?}
% ---------------------------------------------------------------

\subsection{Where We Left Off}

In previous lectures we built a solid foundation using \textbf{linear models}:
\begin{itemize}
  \item Linear combinations of features: $\beta_0 + \beta_1 x_1 + \cdots + \beta_d x_d$
  \item The \textbf{sigmoid} activation $\sigma(x) \to (0,1)$ for producing probabilities
  \item \textbf{Softmax} for multi-class probability distributions
  \item Loss functions: Mean Squared Error (MSE) for regression, Mean Cross-Entropy (MCE) for classification
  \item \textbf{Gradient descent} to find parameters that minimize the loss
\end{itemize}

\begin{notebox}
\textbf{The Key Limitation:} All of these tools only model \emph{linear} relationships.
Real-world data is almost never linearly separable. A fish classifier, for example,
needs curved boundaries to separate Trout, Bass, and Catfish by length and mass ---
a linear model simply cannot do this.
\end{notebox}

\subsection{Universal Function Approximation}

We need a model that, given the right parameter values $\theta$, can approximate \emph{any} function $f(\vec{x}) \to \mathbb{R}$.

Candidate model families include:
\begin{itemize}
  \item \textbf{Multilayer Perceptrons (MLPs)} --- our focus this lecture
  \item Kernel Methods (e.g., SVMs with RBF kernels)
  \item Kolmogorov--Arnold Networks (KANs) --- a recent alternative
\end{itemize}

\begin{takeawaybox}
\textbf{Intuition:} Think of non-linear data as a crumpled piece of paper. A linear model
tries to cut through it with a flat plane --- hopeless. A neural network can \emph{bend}
the space to flatten it first, then apply a linear separator.
\end{takeawaybox}

% ---------------------------------------------------------------
\section{Neural Network Architecture}
% ---------------------------------------------------------------

\subsection{The Basic Building Block}

A neural network stacks \textbf{layers} of computation. Each neuron in a hidden layer:
\begin{enumerate}
  \item Takes a \emph{weighted sum} of all inputs plus a bias
  \item Passes the result through a non-linear \textbf{activation function}
\end{enumerate}

For a simple network with inputs $x_1, x_2$, three hidden neurons, and one output, the hidden activations are:
\[
(a_1 \; a_2 \; a_3) := \sigma\!\left[\,(x_1 \; x_2)\begin{pmatrix} w^1_{1,1} & w^1_{1,2} & w^1_{1,3} \\ w^1_{2,1} & w^1_{2,2} & w^1_{2,3} \end{pmatrix} + (b^1_1 \; b^1_2 \; b^1_3)\,\right]
\]
The output is then:
\[
\hat{y} := (a_1 \; a_2 \; a_3)\begin{pmatrix} w^2_{1,1} \\ w^2_{2,1} \\ w^2_{3,1} \end{pmatrix} + b^2_1
\]
The loss for a single example (regression) is:
\[
\mathcal{L} := (y - \hat{y})^2
\]

\begin{notebox}
\textbf{Key Insight:} A neural network is simply a \emph{composition of functions}.
Each layer transforms its input, and the chain of transformations can represent
highly complex mappings. The activation function is what introduces non-linearity ---
without it, stacking layers would still only produce a linear model.
\end{notebox}

\subsection{Batch Processing}

In practice, we process many examples simultaneously. For a batch of $n$ examples with $d$ input features and $k$ hidden units, the fully-connected layer computes:

\[
Y_{n \times k} := X_{n \times d} \cdot W_{d \times k} + B_{n \times k}
\]

where each row of $B$ is a copy of the bias vector $\vec{b}$. The MSE loss over a batch is:
\[
\vec{r} := y - \hat{y}, \qquad \mathcal{L} := \frac{\vec{r}^T\vec{r}}{n}
\]

Processing in batches is critical for GPU efficiency --- matrix-matrix multiplication is \emph{much} faster than repeated matrix-vector multiplication.

\subsection{The Universal Approximation Theorem}

\begin{theorem}[Universal Approximation, 1989]
A neural network with a single hidden layer of sufficient width and a sigmoid activation function can approximate any continuous function to arbitrary precision.
\end{theorem}

This result, proven in 1989, is the theoretical backbone of deep learning. It tells us that the architecture is expressive enough --- the challenge is \emph{finding} the right weights.

\begin{takeawaybox}
\textbf{Takeaway:} One hidden layer is \emph{theoretically} sufficient. In practice,
deeper networks (more layers) often learn better representations with fewer total
parameters. Width vs.\ depth is an active research area.
\end{takeawaybox}

% ---------------------------------------------------------------
\section{Activation Functions}
% ---------------------------------------------------------------

Activation functions introduce the non-linearity that makes neural networks powerful.
Each has different properties, and the choice matters.

\subsection{Sigmoid}

\[
\sigma(x) = \frac{1}{1 + e^{-x}}, \qquad \sigma'(x) = \sigma(x)(1 - \sigma(x))
\]

Outputs are in $(0,1)$, making it natural for probabilities. Its derivative has a clean closed form. \textbf{Downside:} saturates for large $|x|$, causing \emph{vanishing gradients} deep in a network.

\subsection{ReLU (Rectified Linear Unit)}

\[
\text{relu}(x) = \max(0, x), \qquad \text{relu}'(x) = \begin{cases} 0 & \text{if } x < 0 \\ 1 & \text{if } x > 0 \end{cases}
\]

ReLU is the \textbf{default choice} in modern deep learning. Its gradient is either 0 or 1, which dramatically reduces the vanishing gradient problem. It is also cheap to compute.

\begin{notebox}
\textbf{Dying ReLU:} If a neuron's pre-activation is always negative, its gradient is
always 0 and it \emph{never} learns. This is called the ``dying ReLU'' problem.
Leaky-ReLU was designed to address this.
\end{notebox}

\subsection{Other Common Activations}

\textbf{Tanh:}
\[
\tanh(x) = \frac{e^x - e^{-x}}{e^x + e^{-x}}, \qquad \tanh'(x) = 1 - \tanh^2(x)
\]
Outputs in $(-1, 1)$; zero-centered (often better than sigmoid for hidden layers).

\textbf{Leaky-ReLU:}
\[
\text{leaky-relu}(x) = \max(\alpha x,\, x), \quad \alpha > 0 \text{ (very small)}
\]
\[
\text{leaky-relu}'(x) = \begin{cases} \alpha & \text{if } x < 0 \\ 1 & \text{if } x > 0 \end{cases}
\]
Prevents dying neurons by allowing a small gradient for negative inputs.

\textbf{Swish:}
\[
\text{swish}(x) = x \cdot \sigma(x), \qquad \text{swish}'(x) = \text{swish}(x) + \sigma(x)(1 - \text{swish}(x))
\]
A smooth, self-gated activation that often outperforms ReLU in deep networks.

% ---------------------------------------------------------------
\section{Training Neural Networks: Backpropagation}
% ---------------------------------------------------------------

\subsection{The Core Idea: Everything is Differentiable}

Neural networks are trained by gradient descent, just like linear models. For every parameter $w_{i,j,k}$ we need:
\[
\frac{\partial \mathcal{L}}{\partial w_{i,j,k}} = ?
\]
Then we update: $w \leftarrow w - \eta \cdot \frac{\partial \mathcal{L}}{\partial w}$.

The key observation is that \textbf{every operation in a neural network is differentiable} (or nearly so --- ReLU is not differentiable at 0, but we handle this in practice). This means we can apply the \textbf{chain rule} to propagate gradients backwards through the network.

\subsection{Forward and Backward Passes}

Training alternates between two passes:
\begin{itemize}
  \item \textbf{Forward pass (inference):} Feed data through the network from input to output, computing and storing intermediate activations. Compute the loss.
  \item \textbf{Backward pass (backpropagation):} Walk backwards through the network, computing gradients of the loss with respect to each layer's parameters using the chain rule.
\end{itemize}

\begin{takeawaybox}
\textbf{Analogy:} The forward pass is like driving a car forward; the backward pass is
like reviewing your route on a GPS to see where you should have turned differently.
The gradient tells each parameter: ``increase me to increase the loss, decrease me to
decrease it.''
\end{takeawaybox}

\subsection{Case Study: MNIST Classification}

MNIST is the standard benchmark for image classification:
\begin{itemize}
  \item 60,000 training images and 10,000 test images of handwritten digits
  \item Each image: $28 \times 28$ grayscale pixels, each an 8-bit unsigned integer
  \item Task: classify into one of 10 classes $\{0, 1, \ldots, 9\}$
\end{itemize}

The homework architecture flattens each image to a vector of $28 \times 28 = 784$ values and feeds it through:
\[
\text{Input}(1000 \times 784)
\xrightarrow{\text{FC}, W_1(784,128), b_1(128)}
Z_1(1000 \times 128)
\xrightarrow{\text{ReLU}}
Z_2(1000 \times 128)
\]
\[
\xrightarrow{\text{FC}, W_2(128,10), b_2(10)}
Z_3(1000 \times 10)
\xrightarrow{\text{Softmax}}
\hat{y}(1000 \times 10)
\xrightarrow{\text{MCE Loss}}
\mathcal{L}
\]

\subsection{The Jacobian and the Chain Rule}

For a function $f : \mathbb{R}^d \to \mathbb{R}^m$ with component functions $[f_1(\vec{x}), \ldots, f_m(\vec{x})]$, the \textbf{Jacobian matrix} is:
\[
J_f(x_1, \ldots, x_d) = \begin{pmatrix}
\frac{\partial f_1}{\partial x_1} & \frac{\partial f_1}{\partial x_2} & \cdots & \frac{\partial f_1}{\partial x_d} \\
\frac{\partial f_2}{\partial x_1} & \frac{\partial f_2}{\partial x_2} & \cdots & \frac{\partial f_2}{\partial x_d} \\
\vdots & & \ddots & \vdots \\
\frac{\partial f_m}{\partial x_1} & \frac{\partial f_m}{\partial x_2} & \cdots & \frac{\partial f_m}{\partial x_d}
\end{pmatrix}
\]

The chain rule in matrix form: if $h : \mathbb{R}^m \to \mathbb{R}$, then:
\[
\nabla[h \circ f] = (J_f)^T (\nabla h)
\]

\begin{notebox}
\textbf{Warning:} When activations are matrices (as in batch processing), gradients
are also matrices, and the ``Jacobian'' is technically a four-dimensional tensor.
However, the structure is sparse (many zeros), which we exploit for efficiency.
\end{notebox}

\subsection{Computing the Gradients: Layer by Layer}

We walk backwards through the MNIST network.

\subsubsection{Step 1 --- Gradient at the Output: $\nabla_{Z3}\mathcal{L}$}

The last block is: $Z_3 \xrightarrow{\text{Softmax}} \hat{y} \xrightarrow{\text{MCE}} \mathcal{L}$.

With $\hat{y}$ as the $n \times 10$ prediction matrix and $y$ as the $n \times 10$ one-hot label matrix:
\[
\boxed{\nabla_{Z3}\mathcal{L} = \frac{1}{n}(\hat{y} - y)}
\]
This is a $1000 \times 10$ matrix --- the same shape as $Z_3$ itself.

\subsubsection{Step 2 --- Gradient w.r.t.\ Weights: $\nabla_W \mathcal{L}$}

Since $Z3_{i,j} = b_j + \sum_{k=1}^{128} Z2_{i,k} W_{k,j}$, treating entries of $W$ as variables:
\[
\frac{\partial Z3_{i,j}}{\partial W_{k,a}} = \begin{cases} Z2_{i,k} & \text{if } a = j \\ 0 & \text{if } a \neq j \end{cases}
\]
Collecting these into the chain rule formula:
\[
\boxed{\nabla_W \mathcal{L} = Z2^T \big[\nabla_{Z3}\mathcal{L}\big]}
\]
Shape check: $(128 \times 1000) \cdot (1000 \times 10) = (128 \times 10)$ --- matches $W$.

\subsubsection{Step 3 --- Gradient w.r.t.\ Biases: $\nabla_{\vec{b}}\mathcal{L}$}

Since $\frac{\partial Z3_{i,j}}{\partial b_a} = 1$ if $a = j$ and $0$ otherwise, summing over the batch:
\[
\boxed{\frac{\partial \mathcal{L}}{\partial b_j} = \sum_{i=1}^{n} \big[\nabla_{Z3}\mathcal{L}\big]_{i,j}}
\]
In other words, $\nabla_{\vec{b}}\mathcal{L}$ is the \emph{column sum} of $\nabla_{Z3}\mathcal{L}$.

\subsubsection{Step 4 --- Gradient w.r.t.\ Layer Input: $\nabla_{Z2}\mathcal{L}$}

This gradient is needed to continue backpropagating into the previous layer:
\[
\frac{\partial Z3_{i,j}}{\partial Z2_{a,k}} = \begin{cases} W_{k,j} & \text{if } a = i \\ 0 & \text{if } a \neq i \end{cases}
\]
\[
\boxed{\nabla_{Z2}\mathcal{L} = \big[\nabla_{Z3}\mathcal{L}\big] W^T}
\]
Shape check: $(1000 \times 10) \cdot (10 \times 128) = (1000 \times 128)$ --- matches $Z_2$.

\subsubsection{Step 5 --- Gradient Through ReLU: $\nabla_{Z1}\mathcal{L}$}

ReLU is applied element-wise: $Z2_{i,j} = \max(0, Z1_{i,j})$, so:
\[
\frac{\partial Z2_{i,j}}{\partial Z1_{i,j}} = \begin{cases} 1 & \text{if } Z1_{i,j} \geq 0 \\ 0 & \text{if } Z1_{i,j} < 0 \end{cases}
\]
Therefore:
\[
\boxed{\frac{\partial \mathcal{L}}{\partial Z1_{i,j}} = \begin{cases} \nabla_{Z2}\mathcal{L}_{i,j} & \text{if } Z1_{i,j} \geq 0 \\ 0 & \text{if } Z1_{i,j} < 0 \end{cases}}
\]
This is a simple \emph{masking} operation: pass the upstream gradient through wherever ReLU was active, zero it out elsewhere.

% ---------------------------------------------------------------
\section{Training in Practice}
% ---------------------------------------------------------------

\subsection{Interpreting the Loss Curve}

A healthy training run over 60 epochs on MNIST shows loss decreasing steeply at first,
then leveling off near 0. Key sanity checks:

\begin{itemize}
  \item \textbf{Loss should decrease.} If it increases or overflows your floats, you likely have a bug in your gradient computation.
  \item \textbf{Loss and accuracy should correlate.} Decreasing cross-entropy loss should correspond to improving classification accuracy.
  \item \textbf{Test accuracy matters.} If training loss goes down but test accuracy does not improve, you are overfitting.
\end{itemize}

\begin{notebox}
\textbf{Debugging Tip:} Float overflow (NaN loss) is almost always caused by an incorrect
gradient --- often a sign error or missing the $\frac{1}{n}$ normalization factor. Check your
$\nabla_{Z3}\mathcal{L}$ computation first.
\end{notebox}

% ---------------------------------------------------------------
\section{Summary and Key Takeaways}
% ---------------------------------------------------------------

\begin{takeawaybox}
\textbf{Key Takeaways from Lecture 10}

\begin{enumerate}
  \item \textbf{Neural networks overcome linearity} by composing linear layers with non-linear activation functions, enabling them to model arbitrarily complex functions.

  \item \textbf{Universal Approximation (1989):} A single hidden layer of sufficient width can approximate any continuous function. In practice, depth beats width.

  \item \textbf{The forward pass} computes predictions; the \textbf{backward pass (backprop)} computes gradients of the loss w.r.t.\ every parameter via the chain rule.

  \item \textbf{ReLU is the default activation} for hidden layers: cheap, non-saturating, and empirically effective. Sigmoid/Softmax are reserved for output layers.

  \item \textbf{Batch processing} allows efficient matrix operations. The batch size is a hyperparameter balancing speed and gradient quality.

  \item \textbf{Three gradients per fully-connected layer:}
  \begin{itemize}
    \item $\nabla_W \mathcal{L} = Z_{\text{in}}^T [\nabla_{Z_{\text{out}}} \mathcal{L}]$ (update weights)
    \item $\nabla_{\vec{b}} \mathcal{L} = $ column sum of $\nabla_{Z_{\text{out}}} \mathcal{L}$ (update biases)
    \item $\nabla_{Z_{\text{in}}} \mathcal{L} = [\nabla_{Z_{\text{out}}} \mathcal{L}] W^T$ (propagate upstream)
  \end{itemize}

  \item \textbf{ReLU gradient} is a mask: pass through gradient where input was positive, zero otherwise.
\end{enumerate}
\end{takeawaybox}

\subsection{Further Resources}

\begin{itemize}
  \item \textbf{Reading:} Sections 22.1 -- 22.2.3 of the course textbook
  \item \textbf{Videos:} \textit{Neural Networks} series by 3Blue1Brown (highly recommended for visual intuition on backpropagation)
  \item \textbf{Slides:} by Aaron Hillegass, Georgia Tech
\end{itemize}

\end{document}
